% --- Archon physics formalization source begin ---
% archon:physics
% archon:covers PhyXMiniProblems/problem_phyx_mini_0949.lean
% archon:source-report reports/phyx_mini/problem_phyx_mini_0949.source.json

\chapter{Physics problem phyx\_mini\_0949}
\label{ch:PhyXMiniProblems_problem_phyx_mini_0949}

\paragraph{Problem source.}
\#\# Physical scenario

A uniform bar of length \textbackslash{}( L \textbackslash{}) carries a current \textbackslash{}( I \textbackslash{}) in the direction from point \textbackslash{}( a \textbackslash{}) to point \textbackslash{}( b \textbackslash{}). The bar is in a uniform magnetic field that is directed into the page. Consider the torque about an axis perpendicular to the bar at point \textbackslash{}( a \textbackslash{}) that is due to the force that the magnetic field exerts on the bar.

\#\# Question

A uniform bar of length \textbackslash{}( L \textbackslash{}) carries a current \textbackslash{}( I \textbackslash{}) from end \textbackslash{}( a \textbackslash{}) to \textbackslash{}( b \textbackslash{}), and lies in a uniform magnetic field directed into the page. Find the torque about an axis perpendicular to the bar at point \textbackslash{}( a \textbackslash{}).

\#\# Answer choices

- A: A: IBL\textasciicircum{}2

- B: B: \textbackslash{}frac\{1\}\{2\}IBL\textasciicircum{}2

- C: C: \textbackslash{}frac\{1\}\{4\}IBL\textasciicircum{}2

- D: D: \textbackslash{}frac\{1\}\{3\}IBL\textasciicircum{}2

\#\# Figure caption (auxiliary; use the image as primary evidence)

The image depicts a straight conductor labeled from point \textbackslash{}( a \textbackslash{}) to point \textbackslash{}( b \textbackslash{}) with a current \textbackslash{}( I \textbackslash{}) flowing through it, indicated by a purple arrow. 

The conductor is placed in a uniform magnetic field, represented by blue crosses, which typically denote the field going into the page or screen. 

Point \textbackslash{}( x \textbackslash{}) on the conductor indicates a section with a differential element \textbackslash{}( dx \textbackslash{}). 

Above the magnetic field indicators, a vector \textbackslash{}( \textbackslash{}vec\{B\} \textbackslash{}) is labeled, representing the magnetic field.

The terminals of the conductor at \textbackslash{}( a \textbackslash{}) and \textbackslash{}( b \textbackslash{}) are marked with black dots. The layout suggests examination of the magnetic field's effect on the conductor.

\#\# Dataset metadata

Magnetism; Physical Model Grounding Reasoning, Spatial Relation Reasoning

\paragraph{Recorded answer/context.}
B: B: \textbackslash{}frac\{1\}\{2\}IBL\textasciicircum{}2

\paragraph{Figure/image path.}
/root/proposal\_for\_physic/science-mango/phyx\_mini\_run/phyx\_data/test\_image/949.png

\paragraph{Formalization target.}
create a compiling Lean file with sorry bodies at `PhyXMiniProblems/problem\_phyx\_mini\_0949.lean`. The Lean declarations must preserve the physical quantities, dimensions or dimensional roles, figure labels, governing-law hypotheses, and final relation expressed by this problem.
Use Mathlib/Physlib names found through LeanExplore where available. If a physics API is missing, introduce faithful local abstractions rather than scalar placeholder aliases.

\begin{theorem}[Physics formalization target]
\label{thm:physics:phyx_mini_0949:target}
\lean{PhyXMiniProblems.ProblemPhyXMini0949.problem_phyx_mini_0949}
\uses{def:physics:phyx-mini-0949:phyxminiproblems-problemphyxmini0949-lengthinmeters, def:physics:phyx-mini-0949:phyxminiproblems-problemphyxmini0949-magneticfluxdensityinteslas, def:physics:phyx-mini-0949:phyxminiproblems-problemphyxmini0949-currentinamperes, def:physics:phyx-mini-0949:phyxminiproblems-problemphyxmini0949-torqueinnewtonmeters, def:physics:phyx-mini-0949:phyxminiproblems-problemphyxmini0949-currentcarryingbarsetup, def:physics:phyx-mini-0949:phyxminiproblems-problemphyxmini0949-matchescurrentcarryingbardescription, def:physics:phyx-mini-0949:phyxminiproblems-problemphyxmini0949-matchesprimarycurrentbarfigure, def:physics:phyx-mini-0949:phyxminiproblems-problemphyxmini0949-hasphysicalcurrentbarparameters, def:physics:phyx-mini-0949:phyxminiproblems-problemphyxmini0949-satisfiesstraightbargeometry, def:physics:phyx-mini-0949:phyxminiproblems-problemphyxmini0949-satisfiesuniformintopagemagneticfield, def:physics:phyx-mini-0949:phyxminiproblems-problemphyxmini0949-satisfiesdistributedmagneticforceandtorquelaws}
The assigned autoformalize agent should translate this physics problem into Lean declarations in the covered file, with theorem and lemma proof bodies written as `by sorry`.
\end{theorem}
\begin{proof}
This is an autoformalization task, not a proof task. Produce statements that can later be proved by the physics prover without weakening the physical model.
\end{proof}

% --- Archon declaration topology begin ---
\section{Declaration topology}
\begin{definition}[magnetic Flux Density Dimension]
  \label{def:physics:phyx-mini-0949:phyxminiproblems-problemphyxmini0949-magneticfluxdensitydimension}
  \lean{PhyXMiniProblems.ProblemPhyXMini0949.magneticFluxDensityDimension}
  Magnetic flux density has dimension M T⁻¹ C⁻¹ (tesla)
\end{definition}
\begin{proof}
  This is the declared model definition of magnetic Flux Density Dimension.
\end{proof}
\begin{definition}[electric Current Dimension]
  \label{def:physics:phyx-mini-0949:phyxminiproblems-problemphyxmini0949-electriccurrentdimension}
  \lean{PhyXMiniProblems.ProblemPhyXMini0949.electricCurrentDimension}
  Electric current has dimension C T⁻¹ (ampere)
\end{definition}
\begin{proof}
  This is the declared model definition of electric Current Dimension.
\end{proof}
\begin{definition}[force Per Length Dimension]
  \label{def:physics:phyx-mini-0949:phyxminiproblems-problemphyxmini0949-forceperlengthdimension}
  \lean{PhyXMiniProblems.ProblemPhyXMini0949.forcePerLengthDimension}
  Magnetic force per unit length has dimension M T⁻² (newton/metre)
\end{definition}
\begin{proof}
  This is the declared model definition of force Per Length Dimension.
\end{proof}
\begin{definition}[torque Dimension]
  \label{def:physics:phyx-mini-0949:phyxminiproblems-problemphyxmini0949-torquedimension}
  \lean{PhyXMiniProblems.ProblemPhyXMini0949.torqueDimension}
  Torque has dimension M L² T⁻² (newton metre)
\end{definition}
\begin{proof}
  This is the declared model definition of torque Dimension.
\end{proof}
\begin{definition}[Length Quantity]
  \label{def:physics:phyx-mini-0949:phyxminiproblems-problemphyxmini0949-lengthquantity}
  \lean{PhyXMiniProblems.ProblemPhyXMini0949.LengthQuantity}
  A nonnegative, unit-independent physical length
\end{definition}
\begin{proof}
  This is the declared model definition of Length Quantity.
\end{proof}
\begin{definition}[Magnetic Flux Density Magnitude]
  \label{def:physics:phyx-mini-0949:phyxminiproblems-problemphyxmini0949-magneticfluxdensitymagnitude}
  \lean{PhyXMiniProblems.ProblemPhyXMini0949.MagneticFluxDensityMagnitude}
  \uses{def:physics:phyx-mini-0949:phyxminiproblems-problemphyxmini0949-magneticfluxdensitydimension}
  A nonnegative, unit-independent magnetic-flux-density magnitude
\end{definition}
\begin{proof}
  This is the declared model definition of Magnetic Flux Density Magnitude.
\end{proof}
\begin{definition}[Electric Current Magnitude]
  \label{def:physics:phyx-mini-0949:phyxminiproblems-problemphyxmini0949-electriccurrentmagnitude}
  \lean{PhyXMiniProblems.ProblemPhyXMini0949.ElectricCurrentMagnitude}
  \uses{def:physics:phyx-mini-0949:phyxminiproblems-problemphyxmini0949-electriccurrentdimension}
  A nonnegative, unit-independent conventional-current magnitude
\end{definition}
\begin{proof}
  This is the declared model definition of Electric Current Magnitude.
\end{proof}
\begin{definition}[Force Per Length Magnitude]
  \label{def:physics:phyx-mini-0949:phyxminiproblems-problemphyxmini0949-forceperlengthmagnitude}
  \lean{PhyXMiniProblems.ProblemPhyXMini0949.ForcePerLengthMagnitude}
  \uses{def:physics:phyx-mini-0949:phyxminiproblems-problemphyxmini0949-forceperlengthdimension}
  A nonnegative, unit-independent magnetic force per unit bar length
\end{definition}
\begin{proof}
  This is the declared model definition of Force Per Length Magnitude.
\end{proof}
\begin{definition}[Torque Magnitude]
  \label{def:physics:phyx-mini-0949:phyxminiproblems-problemphyxmini0949-torquemagnitude}
  \lean{PhyXMiniProblems.ProblemPhyXMini0949.TorqueMagnitude}
  \uses{def:physics:phyx-mini-0949:phyxminiproblems-problemphyxmini0949-torquedimension}
  A nonnegative, unit-independent torque magnitude
\end{definition}
\begin{proof}
  This is the declared model definition of Torque Magnitude.
\end{proof}
\begin{definition}[nonnegative SIReadout]
  \label{def:physics:phyx-mini-0949:phyxminiproblems-problemphyxmini0949-nonnegativesireadout}
  \lean{PhyXMiniProblems.ProblemPhyXMini0949.nonnegativeSIReadout}
  Read a nonnegative dimensionful quantity in coherent SI units
\end{definition}
\begin{proof}
  This is the declared model definition of nonnegative SIReadout.
\end{proof}
\begin{definition}[length In Meters]
  \label{def:physics:phyx-mini-0949:phyxminiproblems-problemphyxmini0949-lengthinmeters}
  \lean{PhyXMiniProblems.ProblemPhyXMini0949.lengthInMeters}
  \uses{def:physics:phyx-mini-0949:phyxminiproblems-problemphyxmini0949-lengthquantity, def:physics:phyx-mini-0949:phyxminiproblems-problemphyxmini0949-nonnegativesireadout}
  Read a physical length in metres
\end{definition}
\begin{proof}
  This is the declared model definition of length In Meters.
\end{proof}
\begin{definition}[magnetic Flux Density In Teslas]
  \label{def:physics:phyx-mini-0949:phyxminiproblems-problemphyxmini0949-magneticfluxdensityinteslas}
  \lean{PhyXMiniProblems.ProblemPhyXMini0949.magneticFluxDensityInTeslas}
  \uses{def:physics:phyx-mini-0949:phyxminiproblems-problemphyxmini0949-magneticfluxdensitymagnitude, def:physics:phyx-mini-0949:phyxminiproblems-problemphyxmini0949-nonnegativesireadout}
  Read magnetic flux density in teslas
\end{definition}
\begin{proof}
  This is the declared model definition of magnetic Flux Density In Teslas.
\end{proof}
\begin{definition}[current In Amperes]
  \label{def:physics:phyx-mini-0949:phyxminiproblems-problemphyxmini0949-currentinamperes}
  \lean{PhyXMiniProblems.ProblemPhyXMini0949.currentInAmperes}
  \uses{def:physics:phyx-mini-0949:phyxminiproblems-problemphyxmini0949-electriccurrentmagnitude, def:physics:phyx-mini-0949:phyxminiproblems-problemphyxmini0949-nonnegativesireadout}
  Read conventional current in amperes
\end{definition}
\begin{proof}
  This is the declared model definition of current In Amperes.
\end{proof}
\begin{definition}[force Per Length In Newtons Per Meter]
  \label{def:physics:phyx-mini-0949:phyxminiproblems-problemphyxmini0949-forceperlengthinnewtonspermeter}
  \lean{PhyXMiniProblems.ProblemPhyXMini0949.forcePerLengthInNewtonsPerMeter}
  \uses{def:physics:phyx-mini-0949:phyxminiproblems-problemphyxmini0949-forceperlengthmagnitude, def:physics:phyx-mini-0949:phyxminiproblems-problemphyxmini0949-nonnegativesireadout}
  Read magnetic force per unit length in newtons per metre
\end{definition}
\begin{proof}
  This is the declared model definition of force Per Length In Newtons Per Meter.
\end{proof}
\begin{definition}[torque In Newton Meters]
  \label{def:physics:phyx-mini-0949:phyxminiproblems-problemphyxmini0949-torqueinnewtonmeters}
  \lean{PhyXMiniProblems.ProblemPhyXMini0949.torqueInNewtonMeters}
  \uses{def:physics:phyx-mini-0949:phyxminiproblems-problemphyxmini0949-torquemagnitude, def:physics:phyx-mini-0949:phyxminiproblems-problemphyxmini0949-nonnegativesireadout}
  Read torque magnitude in newton metres
\end{definition}
\begin{proof}
  This is the declared model definition of torque In Newton Meters.
\end{proof}
\begin{definition}[Bar Endpoint]
  \label{def:physics:phyx-mini-0949:phyxminiproblems-problemphyxmini0949-barendpoint}
  \lean{PhyXMiniProblems.ProblemPhyXMini0949.BarEndpoint}
  The two endpoints explicitly labelled in the supplied figure
\end{definition}
\begin{proof}
  This is the declared model definition of Bar Endpoint.
\end{proof}
\begin{definition}[Bar Current Direction]
  \label{def:physics:phyx-mini-0949:phyxminiproblems-problemphyxmini0949-barcurrentdirection}
  \lean{PhyXMiniProblems.ProblemPhyXMini0949.BarCurrentDirection}
  Conventional-current direction along the bar
\end{definition}
\begin{proof}
  This is the declared model definition of Bar Current Direction.
\end{proof}
\begin{definition}[Page Direction]
  \label{def:physics:phyx-mini-0949:phyxminiproblems-problemphyxmini0949-pagedirection}
  \lean{PhyXMiniProblems.ProblemPhyXMini0949.PageDirection}
  Qualitative directions in the plane of, or normal to, the page
\end{definition}
\begin{proof}
  This is the declared model definition of Page Direction.
\end{proof}
\begin{definition}[Page Placement]
  \label{def:physics:phyx-mini-0949:phyxminiproblems-problemphyxmini0949-pageplacement}
  \lean{PhyXMiniProblems.ProblemPhyXMini0949.PagePlacement}
  Relative placement of an endpoint in the supplied diagonal-bar drawing
\end{definition}
\begin{proof}
  This is the declared model definition of Page Placement.
\end{proof}
\begin{definition}[Axis Relation To Bar]
  \label{def:physics:phyx-mini-0949:phyxminiproblems-problemphyxmini0949-axisrelationtobar}
  \lean{PhyXMiniProblems.ProblemPhyXMini0949.AxisRelationToBar}
  Relation of the requested rotation axis to the straight bar
\end{definition}
\begin{proof}
  This is the declared model definition of Axis Relation To Bar.
\end{proof}
\begin{definition}[Bar Model]
  \label{def:physics:phyx-mini-0949:phyxminiproblems-problemphyxmini0949-barmodel}
  \lean{PhyXMiniProblems.ProblemPhyXMini0949.BarModel}
  Material and shape idealization stated for the bar
\end{definition}
\begin{proof}
  This is the declared model definition of Bar Model.
\end{proof}
\begin{definition}[Figure Quantity Label]
  \label{def:physics:phyx-mini-0949:phyxminiproblems-problemphyxmini0949-figurequantitylabel}
  \lean{PhyXMiniProblems.ProblemPhyXMini0949.FigureQuantityLabel}
  Symbolic labels visible in image 949.png
\end{definition}
\begin{proof}
  This is the declared model definition of Figure Quantity Label.
\end{proof}
\begin{definition}[expected Printed Symbol]
  \label{def:physics:phyx-mini-0949:phyxminiproblems-problemphyxmini0949-expectedprintedsymbol}
  \lean{PhyXMiniProblems.ProblemPhyXMini0949.expectedPrintedSymbol}
  \uses{def:physics:phyx-mini-0949:phyxminiproblems-problemphyxmini0949-figurequantitylabel}
  Printed symbol expected for each labelled figure quantity
\end{definition}
\begin{proof}
  This is the declared model definition of expected Printed Symbol.
\end{proof}
\begin{definition}[Current Bar Figure]
  \label{def:physics:phyx-mini-0949:phyxminiproblems-problemphyxmini0949-currentbarfigure}
  \lean{PhyXMiniProblems.ProblemPhyXMini0949.CurrentBarFigure}
  \uses{def:physics:phyx-mini-0949:phyxminiproblems-problemphyxmini0949-barendpoint, def:physics:phyx-mini-0949:phyxminiproblems-problemphyxmini0949-pagedirection, def:physics:phyx-mini-0949:phyxminiproblems-problemphyxmini0949-pageplacement, def:physics:phyx-mini-0949:phyxminiproblems-problemphyxmini0949-figurequantitylabel}
  Literal qualitative content of the supplied raster
\end{definition}
\begin{proof}
  This is the declared model definition of Current Bar Figure.
\end{proof}
\begin{definition}[Current Carrying Bar Setup]
  \label{def:physics:phyx-mini-0949:phyxminiproblems-problemphyxmini0949-currentcarryingbarsetup}
  \lean{PhyXMiniProblems.ProblemPhyXMini0949.CurrentCarryingBarSetup}
  \uses{def:physics:phyx-mini-0949:phyxminiproblems-problemphyxmini0949-lengthquantity, def:physics:phyx-mini-0949:phyxminiproblems-problemphyxmini0949-magneticfluxdensitymagnitude, def:physics:phyx-mini-0949:phyxminiproblems-problemphyxmini0949-electriccurrentmagnitude, def:physics:phyx-mini-0949:phyxminiproblems-problemphyxmini0949-forceperlengthmagnitude, def:physics:phyx-mini-0949:phyxminiproblems-problemphyxmini0949-torquemagnitude, def:physics:phyx-mini-0949:phyxminiproblems-problemphyxmini0949-barendpoint, def:physics:phyx-mini-0949:phyxminiproblems-problemphyxmini0949-barcurrentdirection, def:physics:phyx-mini-0949:phyxminiproblems-problemphyxmini0949-axisrelationtobar, def:physics:phyx-mini-0949:phyxminiproblems-problemphyxmini0949-barmodel, def:physics:phyx-mini-0949:phyxminiproblems-problemphyxmini0949-currentbarfigure}
  Independent physical quantities and spatial data for the bar. The argument of barPointAtMeterCoordinate and of magneticForcePerLengthAtMeter is the coherent-SI coordinate x measured from endpoint a. The requested torque is stored as an independent observable, not defined from an answer choice
\end{definition}
\begin{proof}
  This is the declared model definition of Current Carrying Bar Setup.
\end{proof}
\begin{definition}[page Normal Index]
  \label{def:physics:phyx-mini-0949:phyxminiproblems-problemphyxmini0949-pagenormalindex}
  \lean{PhyXMiniProblems.ProblemPhyXMini0949.pageNormalIndex}
  The page-normal component of a three-dimensional Physlib field vector
\end{definition}
\begin{proof}
  This is the declared model definition of page Normal Index.
\end{proof}
\begin{definition}[Matches Current Carrying Bar Description]
  \label{def:physics:phyx-mini-0949:phyxminiproblems-problemphyxmini0949-matchescurrentcarryingbardescription}
  \lean{PhyXMiniProblems.ProblemPhyXMini0949.MatchesCurrentCarryingBarDescription}
  \uses{def:physics:phyx-mini-0949:phyxminiproblems-problemphyxmini0949-currentcarryingbarsetup}
  Qualitative assumptions stated in the problem prose
\end{definition}
\begin{proof}
  This is the declared model definition of Matches Current Carrying Bar Description.
\end{proof}
\begin{definition}[Matches Primary Current Bar Figure]
  \label{def:physics:phyx-mini-0949:phyxminiproblems-problemphyxmini0949-matchesprimarycurrentbarfigure}
  \lean{PhyXMiniProblems.ProblemPhyXMini0949.MatchesPrimaryCurrentBarFigure}
  \uses{def:physics:phyx-mini-0949:phyxminiproblems-problemphyxmini0949-expectedprintedsymbol, def:physics:phyx-mini-0949:phyxminiproblems-problemphyxmini0949-currentcarryingbarsetup}
  Primary-raster evidence: all printed labels, endpoint placements, uniform crosses, and the current, differential-element, and magnetic-force arrows. No torque value or answer expression occurs here
\end{definition}
\begin{proof}
  This is the declared model definition of Matches Primary Current Bar Figure.
\end{proof}
\begin{definition}[Has Physical Current Bar Parameters]
  \label{def:physics:phyx-mini-0949:phyxminiproblems-problemphyxmini0949-hasphysicalcurrentbarparameters}
  \lean{PhyXMiniProblems.ProblemPhyXMini0949.HasPhysicalCurrentBarParameters}
  \uses{def:physics:phyx-mini-0949:phyxminiproblems-problemphyxmini0949-lengthinmeters, def:physics:phyx-mini-0949:phyxminiproblems-problemphyxmini0949-magneticfluxdensityinteslas, def:physics:phyx-mini-0949:phyxminiproblems-problemphyxmini0949-currentinamperes, def:physics:phyx-mini-0949:phyxminiproblems-problemphyxmini0949-currentcarryingbarsetup}
  Positivity and non-vacuity conditions for the depicted apparatus
\end{definition}
\begin{proof}
  This is the declared model definition of Has Physical Current Bar Parameters.
\end{proof}
\begin{definition}[Satisfies Straight Bar Geometry]
  \label{def:physics:phyx-mini-0949:phyxminiproblems-problemphyxmini0949-satisfiesstraightbargeometry}
  \lean{PhyXMiniProblems.ProblemPhyXMini0949.SatisfiesStraightBarGeometry}
  \uses{def:physics:phyx-mini-0949:phyxminiproblems-problemphyxmini0949-lengthinmeters, def:physics:phyx-mini-0949:phyxminiproblems-problemphyxmini0949-currentcarryingbarsetup}
  The metre parameter starts at endpoint a, ends at endpoint b, and measures distance along the straight bar. Every bar element at the observation time lies in the region occupied by the uniform field
\end{definition}
\begin{proof}
  This is the declared model definition of Satisfies Straight Bar Geometry.
\end{proof}
\begin{definition}[Satisfies Uniform Into Page Magnetic Field]
  \label{def:physics:phyx-mini-0949:phyxminiproblems-problemphyxmini0949-satisfiesuniformintopagemagneticfield}
  \lean{PhyXMiniProblems.ProblemPhyXMini0949.SatisfiesUniformIntoPageMagneticField}
  \uses{def:physics:phyx-mini-0949:phyxminiproblems-problemphyxmini0949-magneticfluxdensityinteslas, def:physics:phyx-mini-0949:phyxminiproblems-problemphyxmini0949-currentcarryingbarsetup, def:physics:phyx-mini-0949:phyxminiproblems-problemphyxmini0949-pagenormalindex}
  The Physlib magnetic field is uniform on the stated region and has only its negative page-normal component, matching the crosses that denote a field into the page. The scalar component is calibrated by the independent dimensionful flux-density magnitude B
\end{definition}
\begin{proof}
  This is the declared model definition of Satisfies Uniform Into Page Magnetic Field.
\end{proof}
\begin{definition}[Satisfies Distributed Magnetic Force And Torque Laws]
  \label{def:physics:phyx-mini-0949:phyxminiproblems-problemphyxmini0949-satisfiesdistributedmagneticforceandtorquelaws}
  \lean{PhyXMiniProblems.ProblemPhyXMini0949.SatisfiesDistributedMagneticForceAndTorqueLaws}
  \uses{def:physics:phyx-mini-0949:phyxminiproblems-problemphyxmini0949-lengthinmeters, def:physics:phyx-mini-0949:phyxminiproblems-problemphyxmini0949-magneticfluxdensityinteslas, def:physics:phyx-mini-0949:phyxminiproblems-problemphyxmini0949-currentinamperes, def:physics:phyx-mini-0949:phyxminiproblems-problemphyxmini0949-forceperlengthinnewtonspermeter, def:physics:phyx-mini-0949:phyxminiproblems-problemphyxmini0949-torqueinnewtonmeters, def:physics:phyx-mini-0949:phyxminiproblems-problemphyxmini0949-currentcarryingbarsetup}
  Governing laws for a current element perpendicular to a uniform magnetic field: the Lorentz-force magnitude per unit length is dF/dx = I B everywhere on the bar; and torque about endpoint a is the first moment of that distributed force, integral\_0\textasciicircum{}L x (dF/dx) dx. The second law deliberately retains the unevaluated interval integral, so it does not assume the requested factor L\textasciicircum{}2 / 2
\end{definition}
\begin{proof}
  This is the declared model definition of Satisfies Distributed Magnetic Force And Torque Laws.
\end{proof}
\begin{definition}[Answer Choice]
  \label{def:physics:phyx-mini-0949:phyxminiproblems-problemphyxmini0949-answerchoice}
  \lean{PhyXMiniProblems.ProblemPhyXMini0949.AnswerChoice}
  Labels of the four answer choices printed with the problem
\end{definition}
\begin{proof}
  This is the declared model definition of Answer Choice.
\end{proof}
\begin{definition}[displayed Expression In SI]
  \label{def:physics:phyx-mini-0949:phyxminiproblems-problemphyxmini0949-answerchoice-displayedexpressioninsi}
  \lean{PhyXMiniProblems.ProblemPhyXMini0949.AnswerChoice.displayedExpressionInSI}
  \uses{def:physics:phyx-mini-0949:phyxminiproblems-problemphyxmini0949-lengthinmeters, def:physics:phyx-mini-0949:phyxminiproblems-problemphyxmini0949-magneticfluxdensityinteslas, def:physics:phyx-mini-0949:phyxminiproblems-problemphyxmini0949-currentinamperes, def:physics:phyx-mini-0949:phyxminiproblems-problemphyxmini0949-currentcarryingbarsetup, def:physics:phyx-mini-0949:phyxminiproblems-problemphyxmini0949-answerchoice}
  Literal transcription of the scalar expressions displayed beside the four answer labels. This metadata does not define the independent torque field
\end{definition}
\begin{proof}
  This is the declared model definition of displayed Expression In SI.
\end{proof}
\begin{definition}[recorded Dataset Answer]
  \label{def:physics:phyx-mini-0949:phyxminiproblems-problemphyxmini0949-recordeddatasetanswer}
  \lean{PhyXMiniProblems.ProblemPhyXMini0949.recordedDatasetAnswer}
  \uses{def:physics:phyx-mini-0949:phyxminiproblems-problemphyxmini0949-answerchoice}
  The answer label recorded in the supplied dataset metadata
\end{definition}
\begin{proof}
  This is the declared model definition of recorded Dataset Answer.
\end{proof}
% --- Archon declaration topology end ---
% --- Archon physics formalization source end ---
